\section{Discussion}
\label{sec:discussion}

\subsection{What Vegas is good at today}

Vegas is, in its current form, a productive tool for the narrow
class of protocols it covers.  A specification of a sealed-bid
auction, an escrow, a voting scheme, or a small game becomes a
file an order of magnitude smaller than the equivalent
hand-written Solidity contract, with the strategic structure
explicit and the front-running mitigation automatic.  The
Gambit, SMT, and witness-record backends provide views that
would otherwise be done by hand and seldom done at all.  The
test harness keeps the views consistent.

The tool is designed for one thing: shortening the loop between
a protocol's specification and the deployed contract that runs
it --- short, mechanical, auditable.  The example library exercises this with thirty-nine
specifications spanning textbook game theory, mechanism design,
and adversarial-pattern smart-contract scenarios.  The pipeline
runs in fractions of a second per example.

\subsection{What Vegas is not yet good at}

Several limitations are worth acknowledging in detail.

\paragraph{Expressive power.}  The language is, by design,
limited to finite types and acyclic protocols.  Loops, dynamic
player counts, ERC-20/ERC-777 token flows, and external contract
calls are all out of scope.  Most real-world protocols use at
least some of these.  We expect to grow the language; the IR is
intended to be robust to those extensions, but the current
language is the bounded fragment that lets every backend stay
exhaustive.

\paragraph{No mechanical equivalence between backends.}  The
backends each read the EventGraph and emit their own artifact.
Equivalence is structural and reviewed in code, not proved.
A regression that ships an inconsistent pair of artifacts would
not be caught by anything stronger than the golden-master test
suite, and that suite checks output text rather than semantic
agreement.  The companion VegasCore (\S\ref{sec:vegascore}) is
the route to a mechanical guarantee here, and the
Solidity-to-EventGraph refinement proof is the central
fellowship-funded continuation of this work.

\paragraph{Fair-play liveness is not statically guaranteed.}  A
contextual guard can become unsatisfiable for fair-playing
participants depending on prior moves.  The language treats this
the same as a strategic quit, which is sometimes the right thing
(if the guard's domain genuinely shrunk) and sometimes too
permissive (if a fair-playing actor was forced out).  A static
satisfiability check over bounded types would address the
common case; it is on the TODO list under the working name
\emph{fair-play liveness}.

\paragraph{Gas costs are not in the model.}  Vegas's Gambit
output ignores gas; equilibrium analyses do not see action
costs.  For some protocols (sealed-bid auctions with small bid
domains) this is the right idealization; for others (high-frequency
games where gas is a strategic variable) it is a real omission.
Bounded per-action upper bounds derivable by tools like
GASTAP~\citep{Albert2020Gastap} could be folded in as
action-cost annotations; this is a planned extension.

\paragraph{MAID is gated, not fixed.}  As detailed in
\S\ref{sec:other-backends}, the MAID encoding silently dropped
context-dependent guards.  The current emitter refuses such
inputs.  A self-only-guard narrowing pass, or a switch to a
constrained-MAID extension, would broaden the MAID backend's
applicability.

\paragraph{The Diamond DAG idea is interesting but
underexplored.}  The Coq/Lean witness-record encoding
(\S\ref{sec:other:diamond}) captures a player's strategy as a
function on prior witnesses.  We have not pushed this far enough
to know what it can prove that ordinary tree-shaped semantics
cannot.  It is, today, an evocative idea: every play is an
inhabitant of a dependent record, and every strategy is a typed
function on histories.

\subsection{Foundations first}

The unifying theme of the limitations is intent rather than
oversight.  Vegas was built with the goal of staying small
enough to describe completely: a finite-graph IR,
finite-type moves, a fixed pool, deterministic guards, no
external calls.
This is a smaller language than its closest neighbours in the
DSL space (Marlowe, Scilla, Move), and far smaller than
general-purpose Solidity.  The smallness is what makes the
multi-backend story coherent.  An EventGraph is exhaustively
explorable; an unbounded loop is not.  A finite move alphabet
yields a finite tree; an unbounded \TT{int} does not.

That smallness is a starting point, not a ceiling.  We expect to
relax pieces of it as we understand which extensions preserve
the analytical properties --- bounded iteration (where the
bound is statically derivable), token flows (with proven
conservation laws), per-round subgames composed from a fixed
library.  The right way to do these extensions is to add them
to the IR carefully and re-derive each backend in turn; the
discipline of the current architecture is what makes this
plausible.

\subsection{Future work}

The most immediate work items are recorded in the project's
\file{TODO.txt}.  Here we mention the larger themes.

\begin{itemize}\itemsep2pt
\item \emph{Verified compilation and verified compiled artifacts.}
  Two complementary routes to mechanical guarantees: a once-and
  for-all proof that the compiler refines the EventGraph
  semantics for any input (the VegasCore-side proof), and a
  per-program certificate emitted alongside each artifact (the
  Diamond DAG route plus a VegasCore-importing backend that
  emits \TT{WFProgram} values with their well-formedness
  obligations discharged by elaboration).  These cover
  different audiences: tool authors care about the first, end
  users about the second.
\item \emph{Bounded extensions.}  Variable deposits, ERC-20
  token flows, and gas-aware liveness (the gas-escrow stage
  policy sketched in the fellowship proposal) are the natural
  next steps.  Each requires care to keep the analytical
  backends faithful to the source.
\item \emph{Operational refinement to the EVM.}  The bail-out
  subgame is the place where the Solidity contract and the
  game model can most plausibly disagree.  A Lean proof that
  the contract refines its EventGraph under the explicit
  execution model (ordering adversary, public mempool, fixed
  timeouts) is the central theorem of the fellowship-funded
  workstream.
\item \emph{Workbench integration.}  The Thrones workbench (an
  external companion project) explores equilibrium analyses
  interactively on the emitted MAIDs and EFGs.  A tighter
  integration --- e.g.\ replaying the analysis as a
  back-annotation on the Vegas source, surfacing dominance and
  best-response findings inline --- is plausible.

\item \emph{A game-relevant program logic.}  Vegas's source
  language presents a protocol whose specification has already
  been stripped of distribution and runtime concerns.  A natural
  question is whether one could lift program logic in the
  Hoare-triple tradition onto this object, with an assertion
  language rich enough to talk about the things a Vegas
  programmer actually cares about: which agent knows what at
  each point in the protocol (epistemic modalities $K_R\,\varphi$
  in the style of dynamic epistemic logic), and what the
  strategic future of the current configuration is worth
  (something like prophecy variables, or open-games-style
  co-utility~\citep{GhaniHedgesWinschel2018CGT} flowing
  contextually backwards from the leaves).  The combination is
  speculative; we do not know what the right primitives are.
  But the building blocks --- a finite event graph with
  explicit visibility, terminal payoff expressions, and
  mechanized strategic semantics in VegasCore --- are unusually
  well-aligned with the ingredients such a logic would need.
  We mark it as a direction worth thinking about, not as a
  planned development.
\end{itemize}

The goal is not a finished tool today.  The goal is a
foundation that grows in a disciplined way, where every
extension has to keep the multi-backend story coherent.  The
present description is intended as the starting point.
